\section{Gospel, teaching, and apostolic office}\label{sec:paul-teaching}

Paul did not write a systematic theology detached from pastoral life. His
arguments arise from baptized Gentiles learning to share one table, grieving
their dead, disputing honor, navigating marriage and celibacy, suing one
another, eating idol-food, testing prophecy, collecting money, and facing
suffering. The resulting theology is coherent around Christ but not reducible
to one slogan.

\subsection{The crucified and risen Messiah}

The gospel Paul says he received and handed on announces that Christ died for
sins, was buried, was raised, and appeared to witnesses (1 Cor 15:1--11). The
cross exposes both human boasting and the wisdom of God; the Resurrection is
not an optional symbol but the ground of proclamation, future hope, and the
transformation of embodied life. Baptism joins believers to Christ's death and
new life; the Eucharistic cup and bread communicate his blood and body and
therefore judge a community that humiliates poor members at the meal.

Paul's favored phrases---``in Christ,'' ``with Christ,'' Christ living in the
believer, and the Spirit dwelling within---describe salvation as participation
and new creation, not only acquittal from a distance. Reconciliation is God's
initiative through Christ and creates an entrusted ministry of
reconciliation. The apostle's own weakness becomes a theater in which power is
recognized as Christ's rather than possessed as Paul's.

\subsection{Grace, faith, works, and law}

Galatians and Romans insist that Gentiles are justified through faith in Jesus
Christ, not by ``works of law.'' The precise scope of works of law remains
debated. Paul directly opposes making
circumcision and Torah observance the boundary Gentiles must cross to belong
to Abraham's family. He also says the law is holy, spiritual, and fulfilled in
love; depicts Gentiles as grafted into Israel's cultivated olive tree; and
expects judgment according to works.

Grace, covenant, election, obedience, and mercy all arise within Israel's
Scriptures. Paul's ``faith apart from works'' stands alongside his demand for
a faith working through love, a Spirit-formed life, care for the poor, and
patient endurance. Catholic teaching receives justification as entirely
initiated by grace and as a real transformation that bears fruit, not a human
achievement or an idle label.

Romans 9--11 is essential. Paul anguishes over his kindred, refuses the claim
that God's word has failed, rejects Gentile boasting over broken branches, and
calls Israel beloved because God's gifts and call are irrevocable. The meaning
of ``all Israel will be saved'' is debated. God's fidelity, mercy, and the
continuing vocation of Israel remain at the center of Paul's argument.

\subsection{One body with many gifts}

Paul addresses local assemblies as God's church and as the body of Christ.
Membership is received in one Spirit through baptism; distinctions of gift
remain, but honor cannot be allocated according to social prestige. Apostles,
prophets, and teachers serve the body; gifts without love become noise. First
Corinthians 13 is therefore not a decorative interlude. It judges charismatic
competition and defines patient, non-possessive love as the more excellent way.

Paul's letters show emerging offices without a single late institutional
chart. Philippians addresses overseers and ministers or deacons; Romans names
Phoebe a deacon; First Corinthians emphasizes charisms and household service;
Paul appoints and sends delegates. The Pastorals give a more developed
ordering of overseers, presbyters, deacons, and widows, but their authorship
and later setting are disputed. Catholic reception rightly hears apostolic
ministry and succession in the whole canon; biography must still date the
forms responsibly.

\subsection{Authority in the form of the cross}

Paul claims genuine authority: he has seen the Lord, was called as apostle,
founded communities, hands on traditions, commands in some matters, and
threatens discipline. He also distinguishes a command of the Lord from his own
judgment, appeals rather than commands Philemon, seeks communal discernment of
prophecy, and says that leaders work with believers rather than dominate their
faith. He refuses a boast that hides weakness.

This tension should not be romanticized. Paul can be controlling, and a letter
arrives without preserving every recipient's voice. Yet his stated norm is the
crucified Christ who did not please himself. Apostolic office is exercised for
building up, tested by the gospel it announces, and made accountable to other
apostles, co-workers, and communities.

\subsection{The Lord's coming and ordinary fidelity}

First Thessalonians expects Christ's coming with vivid nearness and reassures
believers that those who have died will not be excluded. First Corinthians
holds imminent transformation together with resurrection of the dead. Romans
can say salvation is nearer and also plan a western mission. Paul did not give
a date. Second Thessalonians restrains claims that the day has already arrived,
but direct Pauline authorship of that letter is debated.

Expectation produces neither panic nor withdrawal in the undisputed letters.
Paul commands work, care, sexual integrity, generosity, hospitality, prayer,
non-retaliation, submission of life to God, and discernment. Romans 13's call
to subordinate oneself to governing authorities belongs beside Romans 12's
rejection of vengeance and Romans 13's claim that love fulfills the law. The
later political uses of this passage require a separate reception-history
audit and are not treated as established findings here.

\subsection{Bodies, marriage, and sexuality}

Paul regards the body as for the Lord, a member of Christ, and destined for
resurrection. First Corinthians 7 values both marriage and celibacy, requires
mutual conjugal authority, permits separation in some mixed marriages, and
speaks under eschatological pressure. Paul says he is unmarried at the time;
the text does not reveal whether he had never married or was widowed.

Romans 1 and First Corinthians 6 contain condemnations of same-sex acts within
ancient arguments about idolatry, passion, exploitation, holiness, and the
kingdom. The exact reach of Greek sexual terms and their relationship to
modern categories of orientation and identity are debated. Historical context
may not simply erase the canonical moral text; neither may the text be lifted
from Paul's universal indictment of sin, proclamation of mercy, and demand
that Christians renounce judgmental superiority. A history of Catholic moral
reception and modern disagreement requires sources not collated for this
biography and is not claimed here.

\begin{studyframe}
\textbf{Interpretive boundary.} Paul belongs to the Church's canonical
teaching, not only to historical reconstruction. Yet no isolated Pauline
sentence is a complete Catholic doctrine. Read the undisputed letter in its
argument, the letter within the canon, and the canon within the Church's
received teaching; keep later applications visibly later.
\end{studyframe}
